% ============================================================
%  CHAPITRE 6 : Structure d'un code CFD et d'un cas OpenFOAM
% ============================================================
\chapter{Structure d'un code CFD et d'un cas OpenFOAM}
\label{chap:structure_cfd_openfoam}

% SECTION : INTRODUCTION
\section{Introduction}
\label{sec:introduction_structure_cfd_openfoam}

La mécanique des fluides numérique (CFD) permet de simuler des écoulements complexes en résolvant numériquement les équations de conservation. Pour cela, un code CFD suit une structure bien définie, articulée autour de trois étapes principales : le pré-traitement, la résolution et le post-traitement \cite{wendt2009computational}. Chacune de ces phases est indispensable pour garantir la fiabilité des résultats, depuis la préparation du modèle jusqu'à l'analyse des données simulées.

Ce chapitre présente dans un premier temps cette structure générale, puis détaille la manière dont OpenFOAM l'implémente concrètement : son organisation de fichiers, ses solveurs, ses algorithmes de résolution numérique, ainsi que ses outils de visualisation et de post-traitement.

% SECTION : STRUCTURE GÉNÉRALE D'UN CODE CFD
\section{Structure générale d'un code CFD}
\label{sec:structure_generale_code_cfd}

\subsection{Pré-traitement (Pre-processing)}
\label{subsec:preprocessing}

Cette étape consiste à préparer les données nécessaires à la simulation. Cela inclut la définition du domaine géométrique sous forme maillée, le choix des phénomènes physiques à modéliser (comme la turbulence, les écoulements multiphasiques ou la cavitation), ainsi que la spécification des propriétés physiques du fluide et des conditions aux limites adaptées au cas étudié \cite{thompson1999handbook}. Les grandeurs à calculer, telles que la vitesse, la pression et la température, sont associées aux nœuds de chaque cellule du maillage.

La qualité du maillage est cruciale pour la précision des résultats. Les paramètres clés incluent la résolution spatiale, la qualité des cellules (skewness, orthogonalité) et le raffinement local dans les zones à forts gradients \cite{knupp1993fundamentals}.

\subsection{Résolution (Processing)}
\label{subsec:processing}

Le rôle du solveur est d'appliquer les algorithmes permettant \cite{patankar1980numerical} :
\begin{itemize}
    \item D'intégrer les équations de conservation sur chaque volume de contrôle.
    \item De transformer ces équations intégrales en un système algébrique.
    \item D'en estimer la solution à l'aide de méthodes itératives.
\end{itemize}

La majorité des logiciels de CFD, dont OpenFOAM, utilise la méthode des \textbf{volumes finis}, qui repose sur l'application directe des bilans sur chaque volume de contrôle. Cette approche est largement adoptée en raison de sa rigueur et de son efficacité, ce qui explique son succès auprès des éditeurs de logiciels de simulation numérique \cite{versteeg2007introduction}.

\subsection{Post-traitement (Post-processing)}
\label{subsec:postprocessing}

Le post-traitement vise avant tout à évaluer la qualité de la solution obtenue. Il s'agit notamment de vérifier si cette solution est indépendante de la taille du maillage, du critère de convergence et des schémas numériques utilisés \cite{roache1998verification}. Il est également essentiel de s'assurer que le modèle de turbulence ainsi que les conditions aux limites ont été choisis de manière pertinente, et d'analyser dans quelle mesure les résultats dépendent de ces choix.

En général, le solveur fournit une grande quantité de données brutes, telles que les champs de vitesse (u, v, w), de pression (p), etc., pour chaque point du maillage. Ces données doivent être converties en représentations visuelles afin de faciliter leur analyse et leur compréhension \cite{schroeder2006visualization}.

Les principaux résultats généralement présentés lors du post-traitement sont les suivants \cite{hansen2011visualization} :
\begin{itemize}
    \item La représentation de la géométrie du domaine et des détails du maillage.
    \item L'évolution de la convergence au cours des itérations (graphique des résidus).
    \item Les champs et profils de vitesse (contours, vecteurs, lignes de courant).
    \item Les tracés des distributions de pression.
    \item Les animations et visualisations dynamiques des écoulements simulés.
\end{itemize}

% SECTION : STRUCTURE GÉNÉRALE D'UN CAS OPENFOAM
\section{Structure générale d'un cas OpenFOAM}
\label{sec:structure_openfoam}

OpenFOAM adopte une structure de fichiers standardisée, organisée de manière logique pour séparer les différents aspects du cas \cite{openfoam2022guide}. Cette hiérarchie de fichiers facilite la lisibilité, la modification et l'automatisation des cas de calcul. Un cas OpenFOAM comprend trois répertoires principaux : \textbf{``0''}, \textbf{``constant''} et \textbf{``system''} \cite{greenshields2020openfoam}.

% SECTION : LES SOLVEURS DISPONIBLES DANS OPENFOAM
\section{Les solveurs disponibles dans OpenFOAM}
\label{sec:solveurs_openfoam}

OpenFOAM propose une vaste bibliothèque de solveurs classés par type d'écoulement \cite{jasak1996error} :

\begin{table}[H]
    \centering
    \caption{Classification des solveurs dans OpenFOAM}
    \label{tab:solveurs_openfoam}
    \resizebox{\textwidth}{!}{%
    \begin{tabular}{|l|l|l|}
        \hline
        \textbf{Type d'écoulement} & \textbf{Solveurs} & \textbf{Description} \\
        \hline
        Basique & laplacianFoam & Diffusion thermique dans un solide \\
        \hline
        & potentialFoam & Écoulement potentiel pour initialisation \\
        \hline
        Incompressible & icoFoam & Transitoire, laminaire, isotherme \\
        \hline
        & simpleFoam & Permanent, turbulent (algorithme SIMPLE) \\
        \hline
        & pisoFoam & Transitoire, turbulent (algorithme PISO) \\
        \hline
        & pimpleFoam & Transitoire, grand pas de temps (algorithme PIMPLE) \\
        \hline
        Compressible & rhoCentralFoam & Basé sur Kurganov-Tadmor, ondes de choc \\
        \hline
        & rhoPimpleFoam & Transitoire, laminaire ou turbulent \\
        \hline
        & rhoSimpleFoam & Permanent, RANS \\
        \hline
        Multiphase & \textbf{interFoam} & \textbf{2 fluides incompressibles, méthode VOF} \\
        \hline
        & cavitatingFoam & Transitoire, cavitation \\
        \hline
        & multiphaseInterFoam & n fluides incompressibles \\
        \hline
    \end{tabular}}
\end{table}

Pour notre étude, le solveur \textbf{interFoam} est choisi car il est spécifiquement conçu pour les écoulements diphasiques incompressibles avec suivi d'interface par la méthode VOF \cite{rusche2002computational}, ce qui est parfait pour modéliser la surface libre de l'eau dans l'évacuateur de crues.

% SECTION : LES ALGORITHMES DE RÉSOLUTION NUMÉRIQUE
\section{Les algorithmes de résolution numérique}
\label{sec:algorithmes_resolution}

Les équations de Navier-Stokes établissent une relation linéaire entre la pression et la vitesse. Cependant, la pression ne dispose pas d'une équation propre permettant sa résolution directe. Plusieurs algorithmes sont disponibles pour assurer le couplage pression-vitesse \cite{patankar1972calculation} :

\subsection{L'algorithme SIMPLE}
\label{subsec:simple}

L'algorithme SIMPLE (Semi-Implicit Method for Pressure-Linked Equations) est une méthode itérative pour les écoulements \textbf{stationnaires} \cite{caretto1973two}. Il procède par une séquence d'étapes à chaque itération : estimation d'un champ de vitesse à partir d'un champ de pression estimé, puis résolution d'une équation de correction de pression pour satisfaire la continuité, et enfin mise à jour de la pression et de la vitesse. Le processus est répété jusqu'à convergence.

\subsection{L'algorithme PISO}
\label{subsec:piso}

L'algorithme PISO (Pressure Implicit with Splitting of Operator) est spécifiquement adapté aux écoulements \textbf{transitoires} \cite{issa1986solution}. À chaque pas de temps, il effectue une étape de prédiction de la vitesse, suivie de plusieurs étapes de correction de pression (généralement deux) pour satisfaire la continuité de manière plus précise que SIMPLE.

\subsection{L'algorithme PIMPLE}
\label{subsec:pimple}

L'algorithme PIMPLE combine les deux approches, PISO et SIMPLE. Il s'agit d'une implémentation de l'algorithme PISO avec la possibilité d'utiliser un facteur de sous-relaxation (comme dans SIMPLE) et d'effectuer des itérations externes à l'intérieur d'un même pas de temps. Cela permet d'utiliser des pas de temps plus grands que ceux imposés par la condition de Courant, améliorant ainsi l'efficacité pour les simulations transitoires avec des maillages complexes \cite{holzmann2017mathematics}.

% SECTION : STABILITÉ ET CONVERGENCE DES CALCULS
\section{Stabilité et convergence des calculs}
\label{sec:stabilite_convergence}

\subsection{Définitions}
\label{subsec:definitions}

La \textbf{stabilité} en simulation CFD désigne la capacité d'une méthode numérique à générer des résultats bornés et réalistes, sans amplification incontrôlée des erreurs numériques \cite{courant1928uber}. La \textbf{convergence} correspond à la capacité d'une méthode numérique à fournir une solution qui se rapproche de la solution exacte des équations discrétisées à mesure que le calcul itératif progresse (ou que le maillage est raffiné) \cite{ferziger2002computational}.

\subsection{Paramètres influençant la convergence}
\label{subsec:parametres_convergence}

Les principaux paramètres influençant la convergence, la stabilité et la durée d'une simulation CFD sont \cite{hirsch2007numerical} :
\begin{itemize}
    \item La discrétisation spatiale (taille et qualité du maillage).
    \item La discrétisation temporelle (pas de temps $\Delta t$).
    \item Le choix du modèle de turbulence et de ses paramètres.
    \item La méthode de résolution et le type de solveur (SIMPLE, PISO, etc.).
    \item Les conditions aux limites et initiales.
    \item Les propriétés physiques du fluide ($\rho$, $\mu$).
\end{itemize}

\subsection{Critères de stabilité}
\label{subsec:criteres_stabilite}

Pour garantir la stabilité d'un calcul CFD explicite, il est essentiel de respecter le critère de Courant-Friedrichs-Lewy (CFL) \cite{leveque2007finite}. Le nombre de Courant est défini par :

\begin{equation}
Co = \frac{|U| \cdot \Delta t}{\Delta x}
\label{eq:courant}
\end{equation}

La condition de stabilité s'exprime par \( Co \leq 1 \), garantissant que l'information (une particule fluide) ne parcourt pas plus d'une cellule à chaque pas de temps \cite{trefethen1996finite}. Pour les schémas implicites, la condition est moins restrictive, mais un Co trop grand peut encore dégrader la précision.

% SECTION : LE SYSTÈME DE FICHIERS D'OPENFOAM
\section{Le système de fichiers d'OpenFOAM}
\label{sec:fichiers_openfoam}

\subsection{Le dossier ``0''}
\label{subsec:dossier_0}

Le dossier ``0'' contient les dictionnaires définissant les \textbf{conditions initiales et aux limites} pour chaque champ (pression \texttt{p}, vitesse \texttt{U}, fraction volumique \texttt{alpha.water}, paramètres de turbulence \texttt{k}, \texttt{omega}, \texttt{nut}) à l'instant initial ($t=0$) \cite{moukalled2016finite}.

Les principales conditions aux limites disponibles dans OpenFOAM sont \cite{openfoam2022boundary} :

\begin{table}[H]
    \centering
    \caption{Conditions aux limites courantes dans OpenFOAM}
    \label{tab:boundary_conditions}
    \resizebox{\textwidth}{!}{%
    \begin{tabular}{|l|l|l|}
        \hline
        \textbf{Type} & \textbf{Description} & \textbf{Application typique} \\
        \hline
        \texttt{fixedValue} & Valeur fixée par l'utilisateur & Vitesse d'entrée (\texttt{U}) \\
        \hline
        \texttt{zeroGradient} & Gradient normal nul & Sortie libre (\texttt{p}) \\
        \hline
        \texttt{fixedGradient} & Gradient normal fixé & Flux de chaleur imposé \\
        \hline
        \texttt{inletOutlet} & fixedValue en entrée, zeroGradient en sortie & Sortie avec risque de reflux (\texttt{U}) \\
        \hline
        \texttt{totalPressure} & Pression totale constante & Entrée avec pression imposée \\
        \hline
        \texttt{pressureInletOutletVelocity} & U évalué à partir de p & Surface libre, sortie \\
        \hline
        \textbf{\texttt{noSlip}} & \textbf{Vitesse nulle à la paroi} & \textbf{Parois solides} \\
        \hline
        \texttt{slip} & Glissement parfait & Symétrie \\
        \hline
        \texttt{kqRWallFunction} & Condition de paroi pour k, q, R & Parois turbulentes \\
        \hline
        \texttt{omegaWallFunction} & Condition de paroi pour $\omega$ & Parois turbulentes \\
        \hline
    \end{tabular}}
\end{table}

\subsection{Le dossier ``constant''}
\label{subsec:dossier_constant}

Le dossier ``constant'' regroupe les paramètres indépendants du temps \cite{greenshields2020openfoam} :
\begin{itemize}
    \item \textbf{Propriétés physiques} : \texttt{transportProperties} (viscosité, masse volumique, tension superficielle pour les écoulements multiphasiques).
    \item \textbf{Modélisation de la turbulence} : \texttt{turbulenceProperties} (choix du modèle RANS/LES, type de modèle \texttt{kOmegaSST}, etc.).
    \item \textbf{Maillage} : sous-dossier \texttt{polyMesh} contenant la description complète du maillage (points, faces, cellules, frontières).
    \item \textbf{Géométries STL} : sous-dossier \texttt{triSurface} pour les fichiers de surface (utile pour \texttt{snappyHexMesh}).
\end{itemize}

\subsection{Le dossier ``system''}
\label{subsec:dossier_system}

Le dossier ``system'' contient les fichiers de contrôle de la simulation \cite{jasak2009openfoam} :
\begin{itemize}
    \item \textbf{\texttt{controlDict}} : paramètres temporels (startTime, endTime, deltaT, writeInterval).
    \item \textbf{\texttt{fvSchemes}} : schémas de discrétisation pour les dérivées spatiales (\texttt{gradSchemes}, \texttt{divSchemes}, \texttt{laplacianSchemes}) et temporelles (\texttt{ddtSchemes}).
    \item \textbf{\texttt{fvSolution}} : solveurs linéaires, critères de convergence et algorithmes de couplage (SIMPLE, PISO, PIMPLE).
    \item \textbf{\texttt{blockMeshDict}} : génération du maillage structuré de base.
    \item \textbf{\texttt{snappyHexMeshDict}} : configuration du raffinement autour des géométries complexes (STL) et ajout de couches limites.
    \item \textbf{\texttt{decomposeParDict}} : configuration du calcul parallèle (décomposition du domaine).
\end{itemize}

% SECTION : OUTIL DE VISUALISATION ET DE POST-TRAITEMENT (PARAVIEW)
\section{Outil de visualisation et de post-traitement (ParaView)}
\label{sec:paraview}

\textbf{ParaView} constitue l'outil principal de visualisation et de post-traitement des résultats dans OpenFOAM \cite{ahrens2005paraview}. C'est un logiciel open source puissant qui s'appuie sur un module de lecture intégré à OpenFOAM (\texttt{foamToVTK} ou \texttt{paraFoam}). Pour le lancer, il suffit d'exécuter la commande \texttt{paraFoam} depuis le répertoire du cas concerné.

Cet outil très puissant permet d'afficher divers types de données tels que des champs scalaires (pression, température), des vecteurs (vitesse), des lignes de courant ou encore des surfaces de niveau (pour la surface libre avec \texttt{alpha.water}) \cite{ayachit2015paraview}. Il facilite la création d'animations pour visualiser l'évolution des phénomènes simulés.

% SECTION : MODÉLISATION DE LA SURFACE LIBRE - MÉTHODE VOF
\section{Modélisation de la surface libre - Méthode VOF}
\label{sec:vof}

La méthode \textbf{VOF (Volume of Fluid)} est la plus répandue pour suivre l'interface entre deux fluides non miscibles (par exemple, l'eau et l'air) \cite{hirt1981volume}. Elle utilise un champ scalaire \(\alpha\) (fraction volumique) défini par :

\begin{equation}
\alpha = \frac{V_{eau}}{V_{cellule}}
\label{eq:vof_alpha}
\end{equation}

\(\alpha\) varie entre 0 (cellule pleine d'air) et 1 (cellule pleine d'eau). L'interface est située dans les cellules où \(0 < \alpha < 1\). Le transport de \(\alpha\) est régi par l'équation d'advection \cite{rider1998reconstructing} :

\begin{equation}
\frac{\partial \alpha}{\partial t} + \nabla \cdot (\alpha \vec{U}) = 0
\label{eq:vof_transport}
\end{equation}

Les propriétés du fluide (masse volumique \(\rho\) et viscosité \(\mu\)) sont ensuite calculées par une moyenne pondérée \cite{scardovelli1999direct} :

\begin{align}
\rho &= \alpha \rho_{eau} + (1-\alpha)\rho_{air} \label{eq:vof_rho} \\
\mu &= \alpha \mu_{eau} + (1-\alpha)\mu_{air} \label{eq:vof_mu}
\end{align}

Le solveur \texttt{interFoam} d'OpenFOAM implémente cette méthode avec des techniques avancées de compression de l'interface pour maintenir sa netteté.

% SECTION : CONCLUSION
\section{Conclusion}
\label{sec:conclusion_structure_cfd_openfoam}

La simulation CFD repose sur trois étapes essentielles : le pré-traitement, la résolution et le post-traitement. Chacune joue un rôle clé dans la qualité et l'interprétation des résultats. Une bonne préparation du modèle, un solveur fiable et une analyse rigoureuse garantissent des simulations précises et exploitables.

Ce chapitre a également permis d'explorer la manière dont OpenFOAM met en œuvre cette structure générale, en abordant les différents solveurs disponibles, les algorithmes de résolution numérique (SIMPLE, PISO, PIMPLE), ainsi que les notions essentielles de stabilité et de convergence des calculs. L'organisation des dossiers du logiciel a également été examinée, mettant en lumière la clarté et la logique qui facilitent la gestion des simulations. La méthode VOF, implémentée dans le solveur \texttt{interFoam}, est l'outil idéal pour notre modélisation de l'écoulement dans l'évacuateur de crues.

Il ressort de cette étude que la richesse des solveurs et des algorithmes proposés par OpenFOAM permet de traiter une large gamme de problématiques en mécanique des fluides numérique. Le respect des critères de stabilité et de convergence reste essentiel pour garantir la qualité et la fiabilité des résultats obtenus.
